\section*{Chapitre \ref{ch:fm:the-bank-markets-division} --- La division marchés de la banque}

\begin{solution}{exo:fm:the-bank-markets-division:1}
$307.6/5\,302.0=5.8~\%$~; l'exigence de 5.0~\% est respectée, soit $0.05\times5\,302.0=265.1$ milliards de dollars de Tier 1.
\end{solution}

\begin{solution}{exo:fm:the-bank-markets-division:2}
$27\,761/149\,500=18.6\%$. La \omterm{def:fm:the-economics-of-a-trading-firm:roe}{rentabilité des fonds propres} (ROE) du segment de la banque utilise le résultat net applicable aux \omterm{def:fm:the-proprietary-market-making-firm:ownfunds}{fonds propres} ordinaires, après les dividendes privilégiés qu'elle affecte
au segment, ce qui le rend inférieur au résultat net.
\end{solution}

\begin{solution}{exo:fm:the-bank-markets-division:3}
Résultat après impôt $1.4\times0.75=1.05$ milliard de dollars. Rapporté aux RWA~: $1.05/(0.13\times15)=1.05/1.95=53.8~\%$. Rapporté au levier~: $1.05/(0.05\times400)=1.05/20=5.2~\%$.
\end{solution}

\begin{solution}{exo:fm:the-bank-markets-division:4}
$0.12\times0.05=0.6~\%$ par an sur l'exposition au levier. Seuil de rentabilité~: $1.4/400=0.35~\%$ (35~points de base).
\end{solution}

\begin{solution}{exo:fm:the-bank-markets-division:5}
Maxima par desk~: $12.5+20.0+10.0+9.1+4.0+6.5=62.1$. Le groupe a besoin de $\max(35.1, 55.5)=55.5$. Les dérivés actions et le crédit sont contraints par les RWA, tandis que le
groupe est contraint par le levier, si bien que leurs \omterm{def:fm:the-proprietary-market-making-firm:ownfunds}{fonds propres} RWA s'ajoutent à une exigence de levier qui les couvre déjà.
\end{solution}

\begin{solution}{exo:fm:the-bank-markets-division:6}
Avec l'allocation contraignante~: services de prime brokerage (13.5~\%), dérivés actions (18.1~\%) et change (28.1~\%). Avec la seule clé RWA~: taux en flux (15.4~\%) et
repo (53.8~\%).
\end{solution}

\begin{solution}{exo:fm:the-bank-markets-division:7}
Le résultat tombe à 5.625~milliards de dollars (contre 7.05~milliards de dollars), les \omterm{def:fm:the-proprietary-market-making-firm:ownfunds}{fonds propres} à 48.1~milliards de dollars, et la rentabilité à 11.7~\%~: l'essentiel du gain de rentabilité issu de la réduction du
repo est perdu dès qu'une part plus grande des services de prime brokerage la suit.
\end{solution}

\begin{solution}{exo:fm:the-bank-markets-division:8}
Le programme suppose que les autres desks peuvent croître de moitié à marges inchangées et que les desks sont indépendants~; les clients qui se financent
auprès de la banque négocient et compensent aussi avec elle. Réduire par étapes, mesurer l'effet sur les revenus de prime brokerage et de flux, et retarifer le repo avec un \omterm{def:fm:the-bank-markets-division:charge}{coût du bilan}
avant de le réduire.
\end{solution}

\begin{solution}{pb:fm:the-bank-markets-division:1}
\begin{enumerate}
\item Voir les \cref{def:fm:the-bank-markets-division:cet1,def:fm:the-bank-markets-division:lr}.
\item Minimum de 4.5~\%, \omterm{def:fm:the-bank-markets-division:buffers}{coussin de fonds propres de stress} (stress capital buffer) de 2.5~\%, \omterm{def:fm:the-bank-markets-division:buffers}{surcharge G-SIB} de 4.5~\% (méthode~2).
\item Tier 1 fondé sur les risques~: $13.0\%\times1\,981.7=257.6$ milliards de dollars~; \omterm{def:fm:the-bank-markets-division:lr}{ratio de levier}~: $5.0\%\times5\,302.0=265.1$ milliards de dollars~; c'est le \omterm{def:fm:the-bank-markets-division:lr}{ratio de levier} qui est contraignant, de peu.
\item Les coussins statiques du SLR renforcé ont été remplacés par 50~\% de la \omterm{def:fm:the-bank-markets-division:buffers}{surcharge G-SIB} de la société de portefeuille selon la méthode~1, plafonnée à 1~\% pour les filiales bancaires.
\item Par les RWA standardisés et la \omterm{def:fm:the-bank-markets-division:buffers}{surcharge G-SIB}, et par une simulation de l'érosion du capital dans un scénario de stress sévère.
\item Voir la \cref{def:fm:the-bank-markets-division:alloc}.
\item 53.8~\%, 5.2~\% et 5.2~\%.
\item Dérivés actions et crédit~: leurs pondérations de risque sont élevées par rapport à leur bilan.
\item Voir la \cref{prop:fm:the-bank-markets-division:max}.
\item 35.1, 55.5 et 62.1~milliards de dollars~; 11.4~\% sur 62.1~milliards de dollars.
\item Le coût des \omterm{def:fm:the-proprietary-market-making-firm:ownfunds}{fonds propres} que la règle de levier immobilise par unité d'exposition au levier~: $0.12\times0.05=0.6\%$ par an.
\item Taux en flux 0.075, repo $-0.75$, services de prime brokerage 0.45, dérivés actions 1.11, change 0.765, crédit 0.405 (en milliards de dollars).
\item 0.35~\% de l'exposition au levier par an.
\item Il réduit le repo à environ 0.11 de sa taille et fait croître les autres desks à 1.5 fois~; résultat de 9.12~milliards de dollars contre 7.05~milliards de dollars.
\item Que les desks peuvent croître sans réduire leurs marges, et qu'ils sont indépendants.
\item Résultat de 6.08~milliards de dollars (contre 7.05), \omterm{def:fm:the-proprietary-market-making-firm:ownfunds}{fonds propres} de 50.1~milliards de dollars (contre 62.1), rentabilité de 12.1~\% (contre 11.4~\%).
\item Leurs \omterm{def:fm:the-bank-markets-division:alloc}{contraintes de capital contraignantes} diffèrent, si bien que la même activité leur coûte des montants de \omterm{def:fm:the-proprietary-market-making-firm:ownfunds}{fonds propres} différents.
\item Quelle part de ses revenus provient de clients qui se financent aussi auprès du desk repo, et quelle part suivrait le financement.
\item 53.8~\% par RWA, 5.2~\% par levier~; \omterm{def:fm:the-bank-markets-division:charge}{coût du bilan} au seuil de rentabilité de 35~points de base d'exposition au levier par an.
\item Elle consomme peu de capital de risque et beaucoup de bilan, et pour cette division c'est le bilan qui contraint. Tarifée au coût des
\omterm{def:fm:the-proprietary-market-making-firm:ownfunds}{fonds propres} immobilisés par la règle de levier, elle ne franchit pas le taux de rendement minimal exigé (hurdle rate).
\end{enumerate}
\end{solution}

\begin{solution}{iq:fm:the-bank-markets-division:1}
Un ratio fondé sur les risques pondère les expositions selon leur risque~; un \omterm{def:fm:the-bank-markets-division:lr}{ratio de levier} ne le fait pas. Le \omterm{def:fm:the-bank-markets-division:lr}{ratio de levier} sert de filet de sécurité contre des pondérations de risque trop faibles ou
manipulées.
\iqlookfor{le rôle de filet de sécurité, et le fait que chacun contraint des activités différentes.}
\end{solution}

\begin{solution}{iq:fm:the-bank-markets-division:2}
Les pensions livrées sont des prêts garantis à court terme, à faible pondération de risque, mais figurent au bilan pour toute leur taille.
\iqlookfor{pondérations de risque face à la taille du bilan.}
\end{solution}

\begin{solution}{iq:fm:the-bank-markets-division:3}
\omterm{def:fm:the-proprietary-market-making-firm:ownfunds}{Fonds propres} détenus $0.05\times200=10$ milliards de dollars~; coût à 12~\% de 1.2~milliard de dollars par an, plus que le milliard de dollars qu'il rapporte~: il détruit de la valeur.
\iqlookfor{\omterm{def:fm:the-proprietary-market-making-firm:ownfunds}{fonds propres} imposés par la règle de levier, multipliés par le hurdle rate.}
\end{solution}

\begin{solution}{iq:fm:the-bank-markets-division:4}
La somme des maxima par desk dépasse le maximum des sommes, sauf si une même contrainte contraint tous les desks
(voir la \cref{prop:fm:the-bank-markets-division:max}).
\iqlookfor{$\sum\max\ge\max\sum$.}
\end{solution}

\begin{solution}{iq:fm:the-bank-markets-division:5}
Des bilans différents font jouer des contraintes différentes~: une banque contrainte par le levier valorise le bilan cher et le risque bon marché, et l'inverse.
\iqlookfor{la \omterm{def:fm:the-bank-markets-division:alloc}{contrainte de capital contraignante} comme position concurrentielle.}
\end{solution}

\begin{solution}{iq:fm:the-bank-markets-division:6}
Des marges qui baissent à mesure qu'un desk croît, les liens entre desks (clients qui se financent et négocient), la contrainte de stress, les limites à la vitesse de croissance des desks,
et le coût de la réduction d'une franchise.
\iqlookfor{revenus concaves, liens de franchise et coûts d'ajustement.}
\end{solution}
